\documentclass[../main.tex]{subfiles}
\begin{document}
\chapter{Legendre Transform}
\section{Definition and Properties}
\begin{definition}
  The \textit{Legendre transform} of a function $f: D(f) \to \R$ is defined as:
  \[
    f^{*}(\vec{p}) = \sup_{\vec{x} \in D(f)}  [\vec{p} \cdot \vec{x} - f(\vec{x})]
  \]
  with $D(f^{*}) = \{\vec{p} \in \R^{n}: \text{the supremum on the right hand side is finite}\}$.
\end{definition}
\begin{lemma}
  The Legendre transform $f^{*}$ as defined above is convex. \label{legendreConvex}
\end{lemma}
\begin{proof}
  We first need to show that $D(f^{*})$ is a convex set.
  Let $\vec{p}, \vec{q} \in D(f^{*})$ and $t \in (0, 1)$ be arbitrary.
  We want to show that the following is finite:
  \begin{align*}
    &\sup_{\vec{x} \in D(f)} [((1 - t)\vec{p} + t\vec{q})\cdot \vec{x} - f(\vec{x})] \\
    &= \sup_{\vec{x} \in D(f)} [(1 - t)(\vec{p}\cdot \vec{x} - f(\vec{x})) + t(\vec{q} \cdot \vec{x} - f(\vec{x}))] \\
    &\leq (1 - t) \sup_{\vec{x} \in D(f)} (\vec{p} \cdot \vec{x} - f(\vec{x})) + t \sup_{\vec{x} \in D(f)} (\vec{q} \cdot \vec{x} - f(\vec{x}))
  \end{align*}
  The right hand side of this is finite as $\vec{p}, \vec{q} \in D(f^{*})$ and so $(1 - t)\vec{p} + t\vec{q} \in D(f^{*})\ \forall t \in (0, 1)$ and so $D(f^{*})$ is a convex set.

  The above inequality also tells us that:
  \[
    f^{*}((1 - t)\vec{p} + t\vec{q}) \leq (1 - t) f^{*}(\vec{p}) + t f^{*}(\vec{q})\ \forall t \in (0, 1)
  \]
  and so $f^{*}$ is a convex function.
\end{proof}
\begin{proposition}
  If $f$ is convex, then so is $F_{\vec{p}}(\vec{x}) = f(\vec{x}) - \vec{p} \cdot \vec{x}$.
\end{proposition}
\begin{proof}
  For arbitrary $\vec{x}, \vec{y} \in D(f)$ and $t \in (0, 1)$ we have:
  \begin{align*}
    F_\vec{p}((1 - t)\vec{x} + t\vec{y}) &= f((1 - t)\vec{x} + t\vec{y}) - \vec{p} \cdot ((1 - t)\vec{x} + t\vec{y}) \\
                                         &\leq (1 - t)f(\vec{x}) + t f(\vec{y}) - \vec{p} \cdot ((1 - t)\vec{x} + t\vec{y}) \\
                                         &= (1 - t)F_\vec{p}(\vec{x}) + tF_{\vec{p}}(\vec{y})
  \end{align*}
  So $F_{\vec{p}}$ is convex.
\end{proof}
\begin{remark}
  Note that if $f$ is \textbf{strictly} convex, then we can replace $\leq \to <$ and so can conclude that $F_{\vec{p}}$ is strictly convex
\end{remark}
\begin{corollary}
  If $f$ is convex and differentiable then any stationary point of $\vec{p} \cdot \vec{x} - f(\vec{x})$  is a \textbf{global maximum} occurring at the point $\vec{x}(\vec{p})$ which is given by finding $\vec{x}$ such that $\nabla f(\vec{x}) = \vec{p}$. \label{legDiff}

  The Legendre transform of $f$ is then:
  \[
    f^{*}(\vec{p}) = \vec{p} \cdot \vec{x}(\vec{p}) - f(\vec{x}(\vec{p}))
  \]
  with the domain $D(f^{*}) = \{\vec{p} \in \R^{n}: \exists \vec{x} \in \R^{n} \text{ such that } \nabla f(\vec{x}) = \vec{p}\}$.
\end{corollary}
\begin{proof}
  Since $f$ is convex, $f(\vec{x}) - \vec{p} \cdot \vec{x}$ is convex and so $\vec{p} \cdot \vec{x} - f(\vec{x})$ is concave.
  Thus any stationary point of $\vec{p} \cdot \vec{x} - f(\vec{x})$ is a global maximum by \cref{globalMin}.

  To find stationary points of $\vec{p} \cdot \vec{x} - f(\vec{x})$:
  \[
    \nabla(\vec{p}\cdot \vec{x} - f(\vec{x})) = \vec{p} - \nabla f(\vec{x}) = \vec{0} \iff \nabla f(\vec{x}) = \vec{p}
  \]
  so the global maximum occurs at any $\vec{x}(\vec{p})$ with $\nabla f(\vec{x}) = \vec{p}$.

  Since $\vec{x}(\vec{p})$ is a global maximiser of $\vec{p} \cdot \vec{x} - f(\vec{x})$, the Legendre transform of $f$ is:
  \begin{align*}
    f^{*}(\vec{p}) &= \sup_{\vec{x} \in D(f)} [\vec{p} \cdot \vec{x} - f(\vec{x})] \\
                   &= \vec{p} \cdot \vec{x}(\vec{p}) - f(\vec{x}(\vec{p}))
  \end{align*}
  The domain of this is as claimed since there is no global maximum if we do not have a point with $\nabla f(\vec{x}) = \vec{p}$.
\end{proof}
\begin{remark}
  If $f$ is \textbf{strictly} convex, then the solution to $\nabla f(\vec{x}) = \vec{p}$ is unique.

  This is because $\vec{p} \cdot \vec{x} - f(\vec{x})$ would be strictly concave and so the existence of two distinct $\vec{x}(\vec{p})$ that maximise this would imply the existence of a point on the chord joining them which has a strictly greater value, which is a contradiction.
\end{remark}
\begin{example}
  \begin{enumerate}
    \item Consider $f(x) = \frac{1}{2}ax^2$.
      For $a > 0$, this is strictly convex.
      To find $x(p)$, consider:
      \[
        \nabla f(x) = f'(x) = p \iff x = \frac{p}{a}
      \]
      So:
      \begin{align*}
        f^{*}(p) &= p \cdot \frac{p}{a} - \frac{1}{2} a \left(\frac{p}{a}\right)^2 \\
                 &= \frac{p^2}{a} - \frac{p^2}{2a} \\
                 &= \frac{p^2}{2a}
      \end{align*}
      with $D(f^{*}) = \R$
    \item Consider $f(v) = - \sqrt{1 - v^2}$ with $D(f) = (-1, 1)$.
      This is a convex function so we try to find $v(p)$:
      \[
        \nabla f(v) = f'(v) = \frac{v}{\sqrt{1 - v^2}} = p \iff v = \frac{p}{\sqrt{1 + p^2}}
      \]
      so:
      \[
        f^{*}(p) = \frac{p^2}{\sqrt{1 + p^2}} + \frac{1}{\sqrt{1 + p^2}} = \sqrt{1 + p^2}
      \]
      with $D(f^{*}) = \R$.
      Note that $z = f^{*}(p)$ is a branch of a hyperbola.
    \item Consider $f(x) = cx$ for $c > 0$ and $D(f) = \R$.
      This is convex but \textbf{not} strictly convex.

      The term inside the supremum of the Legendre transform is:
      \[
        px - f(x) = (p - c)x
      \]
      The supremum of this only exists for $p = c$ and so $D(f^{*}) = \{c\}$ with $f^{*}(c) = 0$.

      We could have also seen this by considering $\nabla f(x) = p \iff p = c$.
  \end{enumerate}
\end{example}
\begin{theorem}
  If $f$ is convex and twice continuously differentiable then $f^{**} = f$, i.e. the Legendre transform of the Legendre transform of a convex function is just the function itself.
\end{theorem}
\begin{proof}
  From \cref{legendreConvex}, we know that $f^{*}(\vec{p})$ is convex and so we can appeal to \cref{legDiff} to find $f^{**}$.
  Noting that $f^{*}$ is a function of $\vec{p}$, we wish to find $\vec{p}(\vec{x})$ obeying $\nabla_{\vec{p}} f^{*}(\vec{p}(\vec{x})) = \vec{x}$.

  We first need to find $\nabla_{\vec{p}}$ of $f^{*}$ in the form:
  \[
    f^{*}(\vec{p}) = \vec{p} \cdot \vec{x}(\vec{p}) - f(\vec{x}(\vec{p}))
  \]
  Using the multivariate chain rule, we have:
  \begin{align*}
    \pderiv{}{p_i}(f^{*}(\vec{p})) &= x_i + p_j \pderiv{x_j}{p_i} - \at{\pderiv{f}{x_j}}{\vec{x} = \vec{x}(\vec{p})} \pderiv{x_j}{p_i} \\
                                   &= x_i + \pderiv{x_j}{p_i} \underbrace{\left(p_j - [\nabla f(\vec{x}(\vec{p}))]_j\right)}_{=0} \text{ as $\nabla f(\vec{x}(\vec{p})) = \vec{p}$} \\
                                   &= x_i
  \end{align*}
  Note that $\pderiv{x_j}{p_i}$ exists as $f$ is twice differentiable.

  So $\nabla_{\vec{p}} f^{*}(\vec{p}(\vec{x})) = \vec{x}$, thus $\vec{p}(\vec{x}(\vec{p})) = \vec{p}$ so $\vec{p}(\vec{x})$ is indeed the inverse of $\vec{x}(\vec{p})$ and so we have:
  \begin{align*}
    f^{**}(\vec{x}) &= \vec{x} \cdot \vec{p}(\vec{x}) - f^{*}(\vec{p}(\vec{x})) \\
                    &= \vec{x} \cdot \vec{p}(\vec{x}) - (\vec{p}(\vec{x}) \cdot \vec{x} - f(\vec{x}(\vec{p}(\vec{x})))) \\
                    &= f(\vec{x})
  \end{align*}
\end{proof}
\begin{remark}
  Since $f^{**} = f$ is a Legendre transform of the convex function $f^{*}$, the condition that $f$ is convex is necessary.

  However, strict convexity is not needed, for example, returning to $f(x) = cx$ for $c > 0$:
  \[
    f^{**}(x) = \sup_{p \in \{c\}} (xp - f^{*}(p)) = \sup_{p \in \{c\}} (xp) = cx = f(x)
  \]
\end{remark}
\section{Application to Thermodynamics}
Consider a physical system consisting of many molecules, say, $\qty{e23}{}$ molecules.
We also consider only \textit{macroscopic} dynamics, i.e. those which involve length scales much larger than the separation between the molecules.

In \textit{thermal equilibrium}, macroscopic physics can be described by just a few quantities:
\begin{itemize}
  \item Total energy, denoted $E$
  \item Volume, denoted $V$
  \item Temperature, denoted $T$
  \item Pressure, denoted $P$
\end{itemize}

We say a system is \textit{isolated} if it is not iterating with anything else, for example, a gas within a vacuum flask.
In equilibrium, the macroscopic physics of an isolated system is fully specified by $E$ and $V$.
This means that there must be an enormous number $\Omega(E, V)$ of \textit{microscopic} configurations that are identical when viewed from a macroscopic viewpoint.
These numbers are often on the order of $\Omega \sim 10^{10^{23}}$.

\begin{definition}[Entropy]
  We define the \textit{entropy} as:
  \[
    S(E, V) = k_{B} \log \Omega(E, V)
  \]
  where $k_B \approx \qty{1.38e-23}{\joule\kelvin^{-1}}$ is \textit{Boltzmann's constant}.
\end{definition}
If we increase $E$ at fixed $V$, then $S$ also increases as there are more ways to ``partition'' up the total energy between the molecules.
Thus $S$ is strictly increasing with $E$ and so we can invert to find $E(S, V)$.
\begin{definition}[Temperature and Pressure]
  We define the \textit{temperature} as: \label{tempPressure}
  \[
    T = \at{\pderiv{E}{S}}{V} > 0
  \]
  and \textit{pressure} as:
  \[
    P = -\at{\pderiv{E}{V}}{S}
  \]
\end{definition}
\begin{remark}
  We can equivalently rewrite this as:
  \begin{equation}
    \d{E} = T \d{S} - P \d{V} \label{thermoRelation}
  \end{equation}
  which is called the \textit{fundamental thermodynamic relation}.
\end{remark}
\subsubsection{Helmholtz Free Energy}
Suppose we forget about the definition of temperature and consider the Legendre transform with a parameter named $T$ of $E(S, V)$ with respect to $S$ at some fixed $V$.
Call this $-F$:
\[
  -F(T, V) = \sup_{S} [TS - E(S, V)]
\]
In Example Sheet 1 Q4, we will show that in a stable equilibrium, $E$ is a convex function with respect to $S$.
Thus we can appeal to \cref{legDiff} and so this supremum is attained at some $S(T, V)$ which satisfies $\at{\pderiv{E}{S}}{V} = T$ with the Legendre transform given by $T \cdot S(T, V) - E(S(T, V), V)$.

We call:
\[
  F(T, V) = E(S(T, V), V) - T \cdot S(T, V)
\]
the \textit{Helmholtz free energy}.
From \cref{thermoRelation} and the product rule for differentials, we see that:
\[
  \d{F} = \d{E} - (T \d{S} + S \d{T}) = -S \d{T} - P \d{V}
\]
and so:
\[
  S = - \at{\pderiv{F}{T}}{V} \text{ and } P = - \at{\pderiv{F}{V}}{T}
\]
So if we know the Helmholtz free energy as a function of temperature and volume, we can use this to determine the entropy and pressure as functions of temperature and pressure.

\begin{remark}[Use]
  \nonexaminable
  The Helmholtz free energy is useful when considering a \textbf{non-isolated} system in equilibrium with an environment at a temperature $T$, for example a gas in an uninsulated box.
\end{remark}

\subsubsection{Enthalpy}
Suppose now that we instead forget the definition of pressure and take the Legendre transform with a parameter $-P$ of $E(S, V)$ with respect to $V$ keeping $S$ fixed.
Call this $-H(S, P)$:
\[
  -H(S, P) = \sup_{V} [-PV - E(S, V)]
\]
Again, one can show that stability implies that $E$ must be convex with respect to $V$.
Thus, recalling that our parameter was $-P$ and not just $P$, this supremum is attained at some $V(S, P)$ which satisfies $P = - \at{\pderiv{E}{V}}{S}$ with the Legendre transform given by:
\[
  H(S, P) = E(S, V(S, P)) + P \cdot V(S, P)
\]
We call $H$ the \textit{enthalpy} of the system.
From \cref{thermoRelation}, we then obtain:
\[
  \d{H} = \d{E} + P \d{V} + V \d{P} = T \d{S} + S \d{P}
\]
and so:
\[
  T = \at{\pderiv{H}{S}}{P} \text{ and } V = \at{\pderiv{H}{P}}{S}
\]
So if we know the enthalpy as a function of the entropy and pressure, we can use this to determine the temperature and volume as functions of entropy and pressure.
\begin{remark}[Use]
  Enthalpy is a quantity which is often useful in chemistry.
\end{remark}

\subsubsection{Gibbs Free Energy}
Finally, we can take the Legendre transform of $E(S, V)$ with respect to both $S$ and $V$ (i.e. with respect to the input vector $(S, V)$) with parameters $T, -P$, call this $-G(T, P)$.
\begin{remark}
  This is the same as doing the Legendre transform of $F$ with respect to $V$ or equivalently $H$ with respect to $S$.
\end{remark}
Using the definition of the Legendre transform:
\[
  -G(T, P) = \sup_{S, V} \left[\begin{pmatrix}
  S \\
  V \\
  \end{pmatrix} \cdot \begin{pmatrix}
  T \\
  -P \\
  \end{pmatrix} - E(S, V)\right] = \sup_{S, V} [TS - PV - E(S, V)]
\]
Again, by stability $E(S, V)$ is a convex function from $\R^2 \to \R$.
Appealing to \cref{legDiff}, this supremum is attained at some $S(T, P)$ and $V(T, P)$ satisfying:
\[
  \nabla_{(S, V)} E = \begin{pmatrix}
  \at{\pderiv{E}{S}}{V} \\
  \at{\pderiv{E}{V}}{S} \\
  \end{pmatrix} = \begin{pmatrix}
  T \\
  -P \\
  \end{pmatrix} \iff
  T = \at{\pderiv{E}{S}}{V} \text{ and } P = - \at{\pderiv{E}{V}}{S}
\]
with Legendre transform given by:
\[
  G(T, P) = E(S(T, P), V(T, P)) - T \cdot S(T, P) + P \cdot V(T, P)
\]
We call $G$ the \textit{Gibbs free energy}.
Using \cref{thermoRelation} again:
\[
  \d{G} = \d{E} - T \d{S} - S \d{T} + P \d{V} + V \d{P} = - S \d{T} + V \d{P}
\]
and so:
\[
  S = - \at{\pderiv{G}{T}}{P} \text{ and } V = \at{\pderiv{G}{P}}{T}
\]
\begin{remark}
  The functions $E, F, G, H$ are known collectively as \textit{thermodynamic potentials}.
\end{remark}
\subsubsection{Maxwell Relations}
Since partial derivatives commute, we see that:
\begin{align*}
  \frac{\partial^2 E}{\partial S \partial V} &= \frac{\partial^2 E}{\partial V \partial S} \\
  \iff \at{\pderiv{}{S}\left[\at{\pderiv{E}{V}}{S}\right]}{V} &= \at{\pderiv{}{V}\left[\at{\pderiv{E}{S}}{V}\right]}{S} \\
  \iff -\at{\pderiv{P}{S}}{V} &= \at{\pderiv{T}{V}}{S} \text{from \cref{tempPressure}}
\end{align*}
We can also obtain similar relations by again considering the mixed second derivatives of the other potentials $F, G, H$.
These are known together as \textit{Maxwell Relations}.
\end{document}
